\section{Related Work}
\label{sec:related}

\subsection{Static PE Malware Learning}

Static malware analysis uses evidence available without executing a suspect
object. Earlier learned detectors combined hand-engineered program attributes
with neural models \cite{saxe2015deep}, while raw-byte approaches learned
representations directly from complete executables \cite{raff2018malware}.
EMBER established a documented PE feature representation and benchmark
\cite{anderson2018ember}; SOREL-20M increased scale and supplied richer metadata
\cite{harang2020sorel}. Surveys describe the diversity of representations and
learners used in malware analysis \cite{ucci2019survey,gibert2020rise}.

Dataset scale alone does not guarantee realistic evaluation. Random splits can
place related families, acquisition sources, time periods, or near-duplicate
programs across train and test. Security ML also faces distribution shift and an
adaptive adversary, making provenance, split semantics, and operating point
central to a credible result \cite{arp2022dos}. Our study uses a small
identity-audited EMBER subset and random outer folds. It improves row-level
traceability but does not substitute for temporal, family-disjoint, or external
validation.

\subsection{Active Learning and Batch Selection}

Pool-based active learning alternates model fitting, query selection, oracle
annotation, and retraining. Uncertainty sampling is widely used because it is
simple and model compatible \cite{cohn1996active,settles2009active}. Its weakness
in batch settings is redundancy: the highest-scoring samples can occupy the same
region. Core-set acquisition instead seeks representation coverage
\cite{sener2018active}; ensemble methods query disagreement
\cite{beluch2018power}; BatchBALD uses joint mutual information
\cite{kirsch2019batchbald}; and BADGE clusters uncertain gradient embeddings
\cite{ash2020deep}.

These baselines are not interchangeable across model classes. BADGE and
BatchBALD were designed around differentiable or Bayesian deep models, whereas
the present learner is a boosted-tree ensemble. A strong tree-compatible study
should nevertheless compare uncertainty plus an explicit diversity constraint
and a calibrated committee. The current artifact includes random, uncertainty,
and explanation-signal ablations, but not those two stronger families. We treat
that omission as an evidence gate before a general superiority claim.

Explanations can support active learning in different ways. Explainable active
learning has studied the effect of local explanations on the annotator
experience \cite{ghai2021explainable}. Our use is algorithmic: TreeSHAP summaries
enter the acquisition score before annotation. This does not imply that an
attribution with large magnitude or balanced signs yields high expected
information gain. It defines a falsifiable heuristic whose value must be
established against matched baselines.

\subsection{Attribution and Fidelity Evaluation}

LIME approximates a local decision surface \cite{ribeiro2016why}; SHAP unifies
additive attributions through Shapley values \cite{lundberg2017unified}; and
TreeSHAP specializes the computation to tree ensembles
\cite{lundberg2020local}. Explanation evaluation is multi-dimensional. Sanity
checks can expose explanations that do not respond to model parameters
\cite{adebayo2018sanity}, and adversarial constructions can manipulate
post-hoc explanations \cite{slack2020fooling}. Security data add constrained and
attacker-controlled inputs \cite{warnecke2020evaluating}.

Deletion tests measure whether intervening on highly ranked coordinates changes
model output. Their interpretation depends on the replacement distribution:
feature removal can create out-of-distribution inputs and confound the ranking
comparison \cite{hooker2019benchmark,yeh2019infidelity}. Training-derived mean
and median baselines avoid treating zero as universally absent, but they do not
preserve PE structure. Equal-count random masks control the nominal intervention
budget but may frequently choose coordinates unused by a tree model. Stronger
controls include used-feature random selection, magnitude matching, group-aware
replacement, and attribution permutation.

\begin{table*}[t]
\centering
\caption{Positioning relative to representative acquisition and attribution-evaluation methods. ``Required extension'' denotes a comparison that is not present in the retained confirmatory report.}
\label{tab:related}
\scriptsize
\setlength{\tabcolsep}{3pt}
\begin{tabular}{p{3.0cm}p{3.6cm}p{2.5cm}p{2.6cm}p{3.4cm}}
\toprule
Study or family & Principal signal & Explicit batch diversity & Explanation role & Relationship to this work \\
\midrule
Cohn \emph{et al.} \cite{cohn1996active} / uncertainty & Model uncertainty or error reduction & Method dependent & None & Measured baseline \\
Core-set \cite{sener2018active} & Representation coverage & Yes & None & Required extension \\
Ensemble disagreement \cite{beluch2018power} & Predictive disagreement & Indirect & None & Required extension \\
BatchBALD \cite{kirsch2019batchbald} & Joint mutual information & Yes & None & Requires tree-compatible adaptation \\
BADGE \cite{ash2020deep} & Uncertain gradient embeddings & Yes & None & Requires model-compatible adaptation \\
XAL \cite{ghai2021explainable} & Queried samples plus local explanations & Not central & Annotator-facing & Different explanation role \\
Security XAI evaluation \cite{warnecke2020evaluating} & Explanation behavior under security constraints & Not applicable & Evaluation & Motivates controlled diagnostics \\
This work & Margin plus TreeSHAP signed-mass summaries & No & Acquisition and post-hoc diagnostic & Measured vs. random/uncertainty; diversity remains open \\
\bottomrule
\end{tabular}
\end{table*}

The contribution is therefore not the broad claim that combining active learning
and explanations is unprecedented. It is the specific signed-mass acquisition
construction and its repeated, claim-bounded evaluation together with a
controlled assessment of the attribution ranking used by the policy.

